\section{Introduction}

User simulators are established tools for training and evaluating task-oriented dialogue systems when repeated interaction with human users would be costly or difficult to reproduce \citep{schatzmann2007agenda,kreyssig2018neural,sekulic2024reliable}. Existing approaches include agenda-based simulation, corpus-based neural models, and LLM-based simulators \citep{schatzmann2007agenda,kreyssig2018neural,sekulic2024reliable}. ConvLab-3 integrates multiple simulator implementations into a modular toolkit for dialogue-system development and evaluation \citep{zhu2023convlab3}.

Most established user-simulation work considers one user interacting with one system \citep{schatzmann2007agenda,kreyssig2018neural,sekulic2024reliable}. Multi-user interaction adds further coordination problems because several participants may be eligible to speak, an utterance may address one person or the whole group, and participants may maintain different private preferences \citep{mao2024muca,ouchi2016addressee,nonomura2025whospeaks}. MUCA summarizes these additional decisions as what to say, when to respond, and whom to address \citep{mao2024muca}.

This project implements configurable user simulators for option-grounded group decisions with two to seven participants. A generated scenario defines a finite set of public alternatives, while each simulator receives an individual persona, private option preferences, and the behavioral parameters engagement, verbosity, directness, and stubbornness. The participants state initial preferences, discuss the alternatives, narrow the decision space, and cast explicit votes. The possible outcomes are unanimous agreement, a majority decision, or an unresolved result.

The implementation separates participant behavior from conversational coordination. Each simulator creates a structured user action that specifies whether it wants to speak, the communicative function of its contribution, the relevant option or addressee, its reason, and any proposed stance change. The environment handles direct-response obligations, resolves simultaneous attempts to take the conversational floor, advances the broad discussion phases, and counts the final votes. A language model realizes the selected structured action as a short natural-language message.

This division is related to turn-taking accounts that distinguish current-speaker selection, participant self-selection, and speaker continuation \citep{sacks1974turntaking}. Recent multi-agent work likewise combines response obligations with autonomous self-selection instead of assigning all turns through a fixed schedule \citep{nonomura2025whospeaks}. In this project, self-selection belongs to the simulators; the environment only coordinates access to the next turn.

The project is limited to small decision-oriented discussions. It does not attempt to model stable social relationships, status hierarchies, deception, emotional dynamics, or overlapping spoken interaction. Broader agent-based group simulations include richer memory, strategic behavior, and environment dynamics \citep{gu2024agentgroupchat}. The narrower scope keeps simulator decisions inspectable and permits controlled evaluation of participation, utterance length, stance movement, grounding, and voting behavior.
